% ============================================================
%  Příprava na maturitu – Mechatronika
%  Kompiluj: pdflatex maturita_mechatronika.tex
%  Kódování: UTF-8
% ============================================================

\documentclass[a4paper,10pt]{article}

% --- Balíčky ---
\usepackage[T1]{fontenc}
\usepackage[utf8]{inputenc}
\usepackage[czech]{babel}
\usepackage[left=2.4cm, right=1.8cm, top=1.8cm, bottom=1.8cm]{geometry}
\usepackage{lmodern}
\usepackage{microtype}
\usepackage{parskip}
\usepackage{titlesec}
\usepackage{enumitem}
\usepackage{xcolor}
\usepackage{hyperref}
\usepackage{tabularx}
\usepackage{booktabs}
\usepackage{array}
\usepackage{tikz}
\usepackage{eso-pic}
\usepackage{mdframed}
\usepackage{amsmath}
\usepackage{amssymb}
\usepackage{pgfplots}
\pgfplotsset{compat=1.18}
\usetikzlibrary{arrows.meta, positioning, shapes.geometric, calc}

% --- Barvy ---
\definecolor{accent}{HTML}{03254E}
\definecolor{stripeblue}{HTML}{568EA3}
\definecolor{medgray}{HTML}{888888}
\definecolor{lightblue}{HTML}{EAF2F8}
\definecolor{lightyellow}{HTML}{FDFBE4}
\definecolor{lightgreen}{HTML}{EAF7EA}

% --- Modrý pruh vlevo ---
\AddToShipoutPictureBG{%
  \begin{tikzpicture}[remember picture, overlay]
    \fill[stripeblue]
      (current page.north west)
      rectangle
      ([xshift=10mm] current page.south west);
    \fill[accent!60]
      ([xshift=10mm] current page.north west)
      rectangle
      ([xshift=12mm] current page.south west);
  \end{tikzpicture}%
}

\hypersetup{colorlinks=true, urlcolor=accent, linkcolor=accent}

\titleformat{\section}
  {\large\bfseries\color{accent}}
  {\thesection.}
  {0.5em}
  {}
  [\color{accent}\titlerule]

\titlespacing{\section}{0pt}{16pt}{6pt}

\newmdenv[
  backgroundcolor=lightblue,
  linecolor=stripeblue,
  linewidth=1pt,
  roundcorner=4pt,
  innerleftmargin=8pt,
  innerrightmargin=8pt,
  innertopmargin=6pt,
  innerbottommargin=6pt,
  skipabove=4pt,
  skipbelow=4pt
]{pojmybox}

\newmdenv[
  backgroundcolor=lightyellow,
  linecolor=accent!40,
  linewidth=1pt,
  roundcorner=4pt,
  innerleftmargin=8pt,
  innerrightmargin=8pt,
  innertopmargin=6pt,
  innerbottommargin=6pt,
  skipabove=4pt,
  skipbelow=8pt
]{vysvetlenibox}

\newmdenv[
  backgroundcolor=lightgreen,
  linecolor=accent!30,
  linewidth=1pt,
  roundcorner=4p
  t,
  innerleftmargin=8pt,
  innerrightmargin=8pt,
  innertopmargin=6pt,
  innerbottommargin=6pt,
  skipabove=4pt,
  skipbelow=8pt
]{vzorcebox}

\pagestyle{plain}

% ============================================================
\begin{document}

% --- Záhlaví ---
\begin{minipage}[t]{0.75\textwidth}
    {\Huge\bfseries\color{accent} Příprava na maturitu}\\[4pt]
    {\large\color{stripeblue} Mechatronika}\\[2pt]
    \textcolor{medgray}{\small Suchánek Lukáš \quad SPŠ Prosek \quad 2026}
\end{minipage}
\hfill
\begin{minipage}[t]{0.22\textwidth}
    \raggedleft
    \begin{tikzpicture}
        \draw[color=stripeblue, rounded corners=4pt, line width=1pt]
            (0,0) rectangle (3,1.2)
            node[midway, align=center, color=accent, font=\small\bfseries]
            {Otázek: 22};
    \end{tikzpicture}
\end{minipage}

\vspace{6pt}
\noindent\color{accent}\rule{\linewidth}{1.5pt}
\vspace{2pt}

% ============================================================
\section{Spojité řízení a regulace -- PID regulátor}

\begin{pojmybox}
\textbf{Klíčové pojmy:}
zpětná vazba, odchylka, P, I, D složka, proporcionální zesílení,
integrační čas, derivační čas, přeregulace, ustálená odchylka,
ladění (Ziegler-Nichols), přenosová funkce regulátoru
\end{pojmybox}

\begin{vysvetlenibox}
\textbf{Princip zpětnovazebního řízení}\\
Regulace s uzavřenou smyčkou porovnává skutečnou hodnotu $y$ s požadovanou
hodnotou $w$ a vypočítává akční zásah $u$ tak, aby se odchylka $e = w - y$
minimalizovala.

\begin{center}
\begin{tikzpicture}[scale=0.9,
  block/.style={rectangle, draw=accent, fill=lightblue,
                minimum width=1.8cm, minimum height=0.7cm,
                text centered, font=\small},
  sum/.style={circle, draw=accent, fill=lightblue,
              minimum size=0.5cm, font=\small}]
  \node[sum]   (sum)  at (1,0)   {$\Sigma$};
  \node[block] (reg)  at (3,0)   {PID};
  \node[block] (akt)  at (5.2,0) {Aktuátor};
  \node[block] (sus)  at (7.4,0) {Soustava};
  \node[block] (sen)  at (7.4,-1.5) {Senzor};
  \node (w) at (-0.5, 0) {$w$};
  \node (y) at (9.2,  0) {$y$};
  \draw[>-] (w)    -- (sum);
  \draw[>-] (sum)  -- (reg) node[midway,above,font=\scriptsize]{$e$};
  \draw[>-] (reg)  -- (akt) node[midway,above,font=\scriptsize]{$u$};
  \draw[>-] (akt)  -- (sus);
  \draw[>-] (sus)  -- (y);
  \draw[>-] (8.3,0) --  (sen);
  \draw[>-] (sen)   -- (sum) node[near end, left, font=\scriptsize]{$-$};
\end{tikzpicture}
\end{center}

\textbf{Typy regulací podle žádané hodnoty:}
\begin{itemize}[noitemsep, leftmargin=1.4em, topsep=4pt]
    \item \textbf{Stabilizace} -- $w = \text{konst.}$; udržování konstantní teploty, otáček
    \item \textbf{Sledování} -- $w = f(t)$; robot sleduje trajektorii, anténa sleduje družici
    \item \textbf{Programová regulace} -- $w$ se mění podle předem daného programu
          (např. pec: ohřev $\to$ výdrž $\to$ chlazení)
\end{itemize}

\textbf{Matematický popis PID regulátoru}

\begin{vzorcebox}
$$u(t) = K_p \left[ e(t) + \frac{1}{T_i}\int_0^t e(\tau)\,d\tau
+ T_d \frac{de(t)}{dt} \right]$$
Přenosová funkce: $G_R(s) = K_p \left(1 + \frac{1}{T_i s} + T_d s\right)$
\end{vzorcebox}

\textbf{Fyzikální význam parametrů:}
\begin{itemize}[noitemsep, leftmargin=1.4em, topsep=4pt]
    \item \textbf{$K_p$} (proporcionální zesílení) -- určuje sílu reakce na chybu
    \item \textbf{$T_i$} (integrační čas) -- doba, za kterou I-složka "dohoní" P-složku
    \item \textbf{$T_d$} (derivační čas) -- jak dopředu se díváme pro predikci
    \item \textbf{Alternativní zápis:} $K_i = K_p/T_i$, $K_d = K_p \cdot T_d$
\end{itemize}

\textbf{Vliv jednotlivých složek}

\begin{center}
\begin{tabularx}{0.95\linewidth}{@{} l X X X @{}}
    \toprule
    \textbf{Složka} & \textbf{Efekt} & \textbf{Zvyšování}
                    & \textbf{Nebezpečí} \\
    \midrule
    P & Rychlá reakce & Rychlejší odezva & Trvalá odchylka, kmity \\
    I & Odstraní trvale odchylku & Rychlejší eliminace & Přeregulace, nestabilita \\
    D & Tlumí kmitání & Méně překmitu & Zesiluje šum \\
    \bottomrule
\end{tabularx}
\end{center}

\textbf{Projevy jednotlivých složek v čase:}
\begin{itemize}[noitemsep, leftmargin=1.4em, topsep=4pt]
    \item \textbf{P-složka} -- reaguje okamžitě na přítomnou chybu;
          chyba $\to$ akce; problém: zůstává trvalá odchylka
    \item \textbf{I-složka} -- reaguje na minulost (akumuluje chybu);
          čím déle trvá chyba, tím větší akce; eliminuje trvalou odchylku
    \item \textbf{D-složka} -- reaguje na budoucnost (predikce);
          když chyba rychle roste, D-složka brzdí; tlumí překmit
\end{itemize}

\textbf{Problémy v praxi:}
\begin{itemize}[noitemsep, leftmargin=1.4em, topsep=4pt]
    \item \textbf{Derivační ráz} -- skok žádané hodnoty $\to$ D-složka "vyšílí";
          řeší se derivace podle výstupu $y$ místo chyby $e$
    \item \textbf{Windup} -- při saturaci se integrátor "nabije";
          řeší se antiwindup (podmíněná integrace, back-calculation)
    \item \textbf{Šum} -- D-složka zesiluje vysokofrekvenční šum;
          řeší se filtrací měřeného signálu
\end{itemize}

\textbf{Metoda Ziegler-Nichols (oscilační)}\\
Jedna z nejznámějších metod ladění PID regulátoru, vyvinutá v roce 1942.
Postup pro určení optimálních parametrů:

\begin{enumerate}[noitemsep, leftmargin=1.4em, topsep=4pt]
    \item Nastavit $T_i = \infty$, $T_d = 0$ (čistý P regulátor)
    \item Zvyšovat $K_p$ až do trvalých nehasících kmitů (hranice stability)
    \item Odečíst kritické zesílení $K_{kr}$ a kritickou periodu $T_{kr}$
    \item Vypočítat parametry dle tabulky:
\end{enumerate}

\begin{center}
\begin{tabular}{@{} l c c c @{}}
    \toprule
    \textbf{Typ} & $K_p$ & $T_i$ & $T_d$ \\
    \midrule
    P   & $0{,}50\,K_{kr}$ & $\infty$ & $0$ \\
    PI  & $0{,}45\,K_{kr}$ & $0{,}83\,T_{kr}$ & $0$ \\
    PID & $0{,}60\,K_{kr}$ & $0{,}50\,T_{kr}$ & $0{,}125\,T_{kr}$ \\
    \bottomrule
\end{tabular}
\end{center}

\textbf{Nevýhody Ziegler-Nichols:}
\begin{itemize}[noitemsep, leftmargin=1.4em, topsep=4pt]
    \item Systém se musí dostat do kmitů -- nevhodné pro některé provozy
    \item Výsledek je často příliš "agresivní" (velké překmity)
    \item Neplatí pro soustavy s velkým dopravním zpožděním
\end{itemize}

\textbf{Další metody ladění:}
\begin{itemize}[noitemsep, leftmargin=1.4em, topsep=4pt]
    \item \textbf{Chua-Broida} -- založená na přechodové charakteristice;
          změříme čas $T_{63\%}$ a $T_{28\%}$, vypočítáme $T_i$, $T_d$
    \item \textbf{Heuris-tické} -- $T_i \approx 4 \cdot T_{63\%}$, $T_d \approx T_{63\%} / 2$
    \item \textbf{Auto-tuning} -- moderní regulátory samy najdou parametry
          (Siemens, ABB, Honeywell mají vestavěné auto-tuning funkce)
    \item \textbf{Optimalizační} -- minimalizace kritéria (ISE, IAE, ITAE)
          pomocí numerických metod
\end{itemize}

\textbf{Kritéria kvality regulace:}
\begin{itemize}[noitemsep, leftmargin=1.4em, topsep=4pt]
    \item \textbf{ISE} (Integral of Squared Error) -- $\int e^2 dt$;
          penalizuje velké chyby
    \item \textbf{IAE} (Integral of Absolute Error) -- $\int |e| dt$;
          rovnoměrná penalizace
    \item \textbf{ITAE} (Integral of Time-weighted AE) -- $\int t \cdot |e| dt$;
          penalizuje chyby, které přetrvávají dlouho
\end{itemize}

\begin{center}
\begin{tikzpicture}
\begin{axis}[
  width=9cm, height=4.5cm,
  xlabel={Čas $t$ [s]},
  ylabel={Výstup $y(t)$},
  xmin=0, xmax=10, ymin=0, ymax=1.5,
  grid=major, grid style={dotted, medgray},
  legend pos=south east,
  legend style={font=\small},
  tick label style={font=\small},
  label style={font=\small},
  title style={font=\small\bfseries},
  title={Srovnání odezvy regulátorů na skok žádané hodnoty}
]
\addplot[accent, thick, domain=0:10, samples=80]
    {1 - exp(-1.5*x)*(cos(deg(2*x)) + 0.5*sin(deg(2*x)))};
\addlegendentry{PID}
\addplot[stripeblue, dashed, thick, domain=0:10, samples=80]
    {1 - exp(-0.8*x)};
\addlegendentry{PI}
\addplot[red!70!black, dotted, thick, domain=0:10, samples=80]
    {1 - exp(-0.5*x) - 0.1};
\addlegendentry{P (trvalá odch.)}
\addplot[medgray, domain=0:10] {1};
\end{axis}
\end{tikzpicture}
\end{center}

\textbf{Příklad z praxe: Regulace teploty v peci}\\
Máme pec s topným tělesem (0--230\,V) a teploměrem (Pt100).
\begin{itemize}[noitemsep, leftmargin=1.4em, topsep=4pt]
    \item \textbf{Krok 1:} Nastavíme čistý P režim ($T_i=\infty$, $T_d=0$)
    \item \textbf{Krok 2:} Zvyšujeme $K_p$ od 1 do 50; při $K_p=42$ začne systém
          kmitat s periodou $T_{kr} = 120$\,s
    \item \textbf{Krok 3:} Pro PID: $K_p = 0{,}6 \cdot 42 = 25{,}2$,
          $T_i = 0{,}5 \cdot 120 = 60$\,s, $T_d = 0{,}125 \cdot 120 = 15$\,s
    \item \textbf{Krok 4:} Ověříme funkčnost, případně doladíme pro menší překmit
\end{itemize}
\end{vysvetlenibox}

% ============================================================

% ============================================================
\end{document}
